% Propietario conservado: /Users/ruben/Documents/New project/output/EXTREMA_MEDIA_RAZON_RECIPROCIDAD_ES_EN_20260918/ARTICULO/ES/source/sections/nucleo.tex
\section{Formal kernel of Triadic Modular Holography}
\label{x_reciprocidad:sec:nucleo}

Triadic Modular Holography, abbreviated HMT, begins with a finite structure of positions and operations. The order of construction matters: the positions receive arithmetic evaluations, the transitions acquire a ternary signature and the dynamics preserves the information necessary to prolong them. Only afterward are the numerical realizations formed. A constant recognized at the end of this process does not replace the rule that produced its region or the history that determines its evaluation.

The name \emph{All Possible Paths}, APP, expresses the role of paths on the support. It is adopted as the name of the construction, with a conceptual reference to the path formulation of quantum mechanics; it does not imply that its arithmetic law is Feynman's path integral. The \emph{Topological Phase Kernel}, TPK,\footnote{The authors also dedicate the initials TPK to Tan Pengkiang. This dedication is independent of the mathematical definition.} designates the composition of selection, transport, memory update and prolongation. The acronyms are retained only after the objects they represent have been defined.

\subsection{APP: positional support and arithmetic evaluations}

Consider the nine interior marks of a graduated line between zero and ten,
\[
 D_9=\{1,2,3,4,5,6,7,8,9\}.
\]
Their horizontal and vertical arrangement produces the eighty-one positions of \(D_9^2\). The interior marks must not be confused with the ten unit intervals of the reference line. This arrangement fixes an alphabet and two directions; it does not yet use a real expansion to decide a trajectory.

The cyclic identification of the positions leads to the support
\[
 X_9=(\ZZ/9\ZZ)^2.
\]
The chart of positive representatives preserves the symbol nine for the zero class. For every positive integer one defines
\begin{equation}
 \rho_9(n)=1+((n-1)\bmod9),\qquad
 q_9(n)=\frac{n-\rho_9(n)}9,
 \qquad n=\rho_9(n)+9q_9(n).
 \label{x_reciprocidad:eq:residuo-cociente}
\end{equation}
The positive digital root is therefore a choice of representative of the residual class. It does not by itself preserve the integer before reduction; the coordinate \(q_9\) provides precisely the missing information.

At each position \((i,j)\), the accumulation of \(i\) repeated \(j\) times produces \(ij\). The two integer evaluations and their projections are
\[
 S(i,j)=i+j,\qquad P(i,j)=ij,\qquad
 \Sigma(i,j)=\rho_9(i+j),\quad \Pi(i,j)=\rho_9(ij).
\]
The symmetry \((i,j)\mapsto(j,i)\) interchanges the directions of accumulation and preserves both evaluations. It is a property of a single support, not a subsequent identification between two independent matrices.

\begin{definition}[Lifted arithmetic evaluation]
The evaluation of APP that preserves the two sheets is
\begin{equation}
 \mathcal A(i,j)=
 \bigl((R^+_{ij},q^+_{ij}),(R^\times_{ij},q^\times_{ij})\bigr)
 =\bigl((\rho_9(i+j),q_9(i+j)),(\rho_9(ij),q_9(ij))\bigr).
 \label{x_reciprocidad:eq:app-levantada}
\end{equation}
The term \emph{sheet} here indicates an evaluation on the same support: additive or multiplicative. The two quotient coordinates belong to the arithmetic state and are not discarded when reducing the residues.
\end{definition}


% BEGIN OWNER /Users/ruben/Documents/New project/output/EXTREMA_MEDIA_RAZON_RECIPROCIDAD_ES_EN_20260918/ARTICULO/ES/source/figures/app.tex
\begin{figure}[tbp]
\centering
\begin{tikzpicture}[x=.53cm,y=.53cm,font=\small]
\fill[black!5] (1.5,-.5) rectangle (2.5,8.5);
\fill[black!5] (4.5,-.5) rectangle (5.5,8.5);
\fill[black!5] (-.5,2.5) rectangle (8.5,3.5);
\fill[black!5] (-.5,5.5) rectangle (8.5,6.5);
\fill[black!12] (7.5,-.5) rectangle (8.5,8.5);
\fill[black!12] (-.5,-.5) rectangle (8.5,.5);
\foreach \i in {1,...,9}{
 \node at (\i-1,9.05) {\(\i\)};
 \node at (-1.05,9-\i) {\(\i\)};
 \foreach \j in {1,...,9}{
  \pgfmathtruncatemacro{\v}{mod(\i*\j-1,9)+1}
  \node at (\j-1,9-\i) {\(\v\)};
 }
}
\foreach \k in {0,...,9}{
 \draw[black!25] (\k-.5,-.5)--(\k-.5,8.5);
 \draw[black!25] (-.5,\k-.5)--(8.5,\k-.5);
}
\draw[line width=.65pt] (2.5,3.5) rectangle (4.5,5.5);
\node[anchor=west,align=left,text width=6.5cm] at (10,7.2) {One support: \(81\) positions.\\[5pt]Two evaluations:\\ \(S(i,j)=i+j\), \(P(i,j)=ij\).\\[5pt]Each row is generated by accumulation:\\ \(P(i,j+1)-P(i,j)=i\).};
\node[anchor=west,align=left,text width=6.5cm] at (10,1.8) {La carta representa \(0\pmod9\) por \(9\).\\[5pt]El bloque señalado es\\ \(B_c=\left(\begin{smallmatrix}7&2\\2&7\end{smallmatrix}\right)\),\\ con autovalores \(9\) y \(5\).};
\end{tikzpicture}
\caption{Tabla multiplicativa de APP en representantes positivos. Las filas y columnas no unitarias se distinguen por trama gris; el marco identifica el bloque diagonal adyacente de diagonales iguales. La tabla dibujada es un observable del toro aritmético; sus bordes de representación no constituyen una frontera topológica intrínseca del toro.}
\label{x_reciprocidad:fig:app}
\end{figure}

% END OWNER


\begin{proposition}[Reconstruction of the evaluations]
Los dos pares de \eqref{x_reciprocidad:eq:app-levantada} determinan exactamente \(i+j\) e \(ij\). Los residuos aislados no determinan esas evaluaciones enteras ni sus cocientes.
\end{proposition}
\begin{proof}
La primera afirmación es \eqref{x_reciprocidad:eq:residuo-cociente} aplicada a cada hoja. Para la segunda, \(10=1+9\) y \(19=1+2\cdot9\) tienen el mismo residuo y distinto cociente. En el soporte bidimensional, las posiciones \((5,5)\) y \((8,2)\) dan la misma suma y el mismo residuo producto, pero productos \(25=7+18\) y \(16=7+9\). La proyección residual identifica datos que la evaluación levantada distingue.
\end{proof}

\subsubsection{Grafo de Cayley, topología toroidal y grado de enrollamiento}

Las traslaciones cardinales están dadas, en coordenadas de fila y columna, por
\[
 \nu_N=(-1,0),\quad\nu_E=(0,1),\quad
 \nu_S=(1,0),\quad\nu_O=(0,-1),\qquad T_d(x)=x+\nu_d.
\]
The first coordinate increases southward and the second eastward. The graph
\[
 \mathscr G_9=\operatorname{Cay}(X_9,\{\nu_N,\nu_E,\nu_S,\nu_O\})
\]
is the two-dimensional cyclic grid: it has eighty-one vertices and four outgoing oriented edges per vertex. It is a Cayley graph for the \emph{additive} structure of the support. Multiplication of residues has zero divisors and does not make all rows of APP into a multiplicative group.

A cardinal word \(w=d_1\cdots d_r\) and an origin \(x_0\) determine
\[
 x_k=x_0+\sum_{j=1}^k\nu_{d_j}\pmod9.
\]
There are \(81\cdot4^r\) oriented paths of length \(r\), including returns and backtracking. When the path closes, its integer displacement before reduction defines
\begin{equation}
 \operatorname{wind}(w)=\frac19
 \bigl(\#S(w)-\#N(w),\#E(w)-\#O(w)\bigr)\in\ZZ^2.
\end{equation}
Reversing the path changes the sign of this vector. Returning to the same residual position therefore does not require the winding to be zero. This is the first example of a difference that will reappear in the phase return: one coordinate may close while another records a nontrivial evolution.

\subsubsection{Transport quotient and cocycle identity}

Let \(s_9:\ZZ/9\ZZ\to D_9\) be the section of positive representatives. The transport quotient associated with additive composition, also called \emph{carry} in positional arithmetic, is
\[
 c_9(a,b)=\frac{s_9(a)+s_9(b)-s_9(a+b)}9.
\]
The compositional character of this memory is expressed by an exact identity.

\begin{lemma}[Cocycle identity]
For \(a,b,c\in\ZZ/9\ZZ\),
\[
 c_9(a,b)+c_9(a+b,c)=c_9(b,c)+c_9(a,b+c).
\]
\end{lemma}
\begin{proof}
Both sums are equal to
\[
 \frac{s_9(a)+s_9(b)+s_9(c)-s_9(a+b+c)}9.
\]
The associativity of the sum of lifted representatives requires precisely this additional term when composing in the quotient.
\end{proof}

With positive representatives, this cocycle is not normalized at the identity: \(c_9(0,a)=1\). If \(s_0\) takes representatives in \(\{0,\ldots,8\}\), its normalized carry \(c_0\) satisfies
\[
 c_9(a,b)=c_0(a,b)+\mathbf1_{a=0}+\mathbf1_{b=0}-\mathbf1_{a+b=0}.
\]
The difference is the change of section; it does not alter the cocycle identity. The subsequent decomposition of the calendar into phase and window will use representatives from zero to eight.

The totals of both evaluations distinguish the visible contribution and the one preserved in the quotient:
\[
 \sum S=810,\quad\sum P=2025,\quad
 \sum\Sigma=405,\quad\sum\Pi=459,
\]
\[
 \sum q^+=45,\qquad\sum q^\times=174.
\]
It follows that
\begin{equation}
 2025-810=(459-405)+9(174-45)=54+9\cdot129.
 \label{x_reciprocidad:eq:defecto-integrado}
\end{equation}
The sum--product difference is not exhausted by the visible defect \(54\). The identity also preserves the quotient defect \(129\). Removing it would change the arithmetic being transported, even if the residual table remained identical.

\subsubsection{Ternary reduction and multiplicative radical}

The ring \(R=\ZZ/9\ZZ\) has unit group \(\{1,2,4,5,7,8\}\) and maximal ideal
\[
 \mathfrak m=(3)=\{0,3,6\},\qquad\mathfrak m^2=0.
\]
Every additive row is a permutation of the nine representatives. A multiplicative row is one exactly when its index is a unit. Rows three and six contain three repeated values; row nine contains only the representative of the zero class. Thus the multiplicative folding already distinguishes, before any continuous geometry, the units from the nilpotent sector.

The map \(x\mapsto3x\) has kernel and image both equal to \(\mathfrak m\). It induces an isomorphism of vector spaces over \(\FF_3\),
\[
 R/\mathfrak m\longrightarrow\mathfrak m,\qquad [x]\longmapsto3x,
\]
but not a ring isomorphism. In particular, the visible sequence \(3,6,9\) admits two different readings: its direct reduction modulo three is \(0,0,0\), whereas extracting the factor three followed by reducing its internal coordinate gives \(1,2,0\). This second operation preserves a coordinate of the ideal; it is not the same projection as directly reducing the original number.

\subsubsection{Latin realization of the additive evaluation}

The difference between the two evaluations can also be expressed through
operators. Let \((e_r)_{r\in\ZZ/9\ZZ}\) be the orthonormal basis of \(\CC^9\).
The assignment \(v_{ab}=e_{a+b}\) produces an orthonormal basis in each row
and column. The projectors \(P_{ab}=v_{ab}v_{ab}^{*}\) satisfy
\[
 \sum_bP_{ab}=I_9=\sum_aP_{ab},\qquad P_{ab}^2=P_{ab}.
\]
Indeed, each translation \(b\mapsto a+b\) permutes the nine residues.
A realization of a Latin square by rank-one projectors is obtained,
in the terminology of quantum Latin squares
\cite{x_reciprocidad:mustovicary,x_reciprocidad:delascuevas2023}. This construction starts from the additive
sheet already defined and makes explicit the map to its operator representation.

The multiplicative evaluation preserves other content: for \(a=3\), the
vectors \(e_{3b}\) repeat the residues \(0,3,6\), and
\[
 \sum_{b\in\ZZ/9\ZZ}|e_{3b}\rangle\langle e_{3b}|
 =3\bigl(|e_0\rangle\langle e_0|+|e_3\rangle\langle e_3|
                    +|e_6\rangle\langle e_6|\bigr).
\]
The multiplicity is recorded in the operator and characterizes the multiplicative
radical. Thus the toroidal support, the Latin property and the
operator resolution of the identity are properties that require
different definitions. The quantum magic squares studied by
De las Cuevas, Drescher and Netzer require positive entries and row
and column sums equal to the identity; their theory also distinguishes
dilations with commutative entries\cite{x_reciprocidad:delascuevas2020}. This reference
provides a precise language for comparing realizations of APP.

\subsection{TRIT: orientation and quadratic regimes}

\begin{definition}[Ternary signature and local lift]
The TRIT signature uses the balanced identification
\[
 \FF_3\longrightarrow\TRIT,\qquad 0\mapsto0,\quad1\mapsto+1,\quad2\mapsto-1.
\]
For \(n\in\ZZ\), its lift preserves \(n=r+3c\), where \(r\in\TRIT\) and \(c\in\ZZ\). In a transition of bounded local amplitude, the three states \((0,-1),(0,0),(0,+1)\) may have the same zero residue and different carries.
\end{definition}

The signature does not replace the trajectory. A residual zero does not mean an absence of orientation, phase or quotient information. Nor does it by itself introduce a physical quantity. Its function is to distinguish algebraically the transport classes that will be realized next.

\subsubsection{Three quadratic regimes and a common exponential}

Once the ternary signature has been obtained, its quadratic realization is defined by
\begin{equation}
 \mathbb A_\tau=\RR[X]/(X^2+\tau),\qquad \tau\in\TRIT.
 \label{x_reciprocidad:eq:algebras-trit}
\end{equation}
If \(J_\tau\) is the class of \(X\), then \(J_\tau^2=-\tau I\). The structure distinguishes the elliptic type \(\tau=+1\), the parabolic type \(\tau=0\) and the hyperbolic type \(\tau=-1\). The real realization does not generate the signature: it expresses the types that oriented reduction has already produced.

\begin{proposition}[Exponentiation of the three types]
In each algebra \eqref{x_reciprocidad:eq:algebras-trit},
\begin{equation}
 \exp(tJ_\tau)=C_\tau(t)I+S_\tau(t)J_\tau,
\end{equation}
where
\[
 C_\tau(t)=\sum_{n\ge0}\frac{(-\tau)^nt^{2n}}{(2n)!},\qquad
 S_\tau(t)=\sum_{n\ge0}\frac{(-\tau)^nt^{2n+1}}{(2n+1)!}.
\]
The three realizations are
\[
 \begin{array}{c|c|c}
 \tau&J_\tau^2&\exp(tJ_\tau)\\\hline
 +1&-I&\cos t\,I+\sin t\,J_\tau\\
 0&0&I+tJ_\tau\\
 -1&I&\cosh t\,I+\sinh t\,J_\tau.
 \end{array}
\]
\end{proposition}
\begin{proof}
Separating the even and odd terms of the exponential series and substituting \(J_\tau^{2n}=(-\tau)^nI\) and \(J_\tau^{2n+1}=(-\tau)^nJ_\tau\) gives the formulas. For \(\tau=0\), all terms of degree at least two vanish. In the other two cases the trigonometric or hyperbolic series are recovered.
\end{proof}

The exponential is common to the three regimes: it is not assigned exclusively to the parabolic branch. The invariants
\[
 C_\tau(t)^2+\tau S_\tau(t)^2=1
\]
follow from \(\exp(tJ_\tau)\exp(-tJ_\tau)=I\). The same algebraic principle preserves unity across realizations that differ in their quadratic type.

In the hyperbolic regime, \(P_\pm=(I\pm J_{-1})/2\) are complementary projectors and
\[
 \exp(tJ_{-1})=e^tP_++e^{-t}P_-.
\]
Growth and contraction appear jointly in two eigendirections. In the elliptic regime, the parameter describes a rotation; in the parabolic regime, a nilpotent translation. The temporal interpretation adopted by HMT associates \(+1\) with the return with memory, \(0\) with the threshold and \(-1\) with bidirectional opening. This interpretation requires an order of parametrization; it does not add a fourth algebraic type.


% BEGIN OWNER /Users/ruben/Documents/New project/output/EXTREMA_MEDIA_RAZON_RECIPROCIDAD_ES_EN_20260918/ARTICULO/ES/source/sections/trit_desarrollo.tex
% Ampliación acumulativa del TRIT; se inserta antes de la figura de regímenes.
% No sustituye ni elimina el desarrollo precedente.
\subsubsection{Balanced division and composition of lifts}
\label{x_reciprocidad:trit-ampl-div-balanceada}

Ternary reduction allows an integer evaluation of APP to be expressed in a
visible coordinate and a prolongation coordinate. Its function does not consist
merely in abbreviating the notation of an integer: it provides a composition
rule that determines how the quotient is modified when evaluations are added or
multiplied. The balanced choice also makes explicit the
compatibility of that rule with sign reversal.

\begin{definition}[Balanced ternary coordinates]
\label{x_reciprocidad:trit-ampl-def-coordenadas}
For \(n\in\mathbb Z\), let
\[
 q_3(n)=\left\lfloor\frac{n+1}{3}\right\rfloor,\qquad
 r_3(n)=n-3q_3(n).
\]
The map
\[
 \mathcal B:\mathbb Z\longrightarrow\TRIT\times\mathbb Z,\qquad
 n\longmapsto(r_3(n),q_3(n))
\]
is called the balanced ternary decomposition. Its inverse is
\(\mathcal L(r,c)=r+3c\).
\end{definition}

Indeed, \(r_3(n)\in\{-1,0,+1\}\). If two pairs represent the same integer,
\(r+3c=s+3d\), then \(r-s\) is a multiple of \(3\) and belongs to the interval
\([-2,2]\); necessarily \(r=s\) and \(c=d\). The representation is therefore
unique. Division can be repeated on \(q_3(n)\), so that each level
preserves the datum necessary to reconstruct the preceding one. The quotient is not a
perturbation of the residue: it constitutes its complementary coordinate.

The signature obtained from an evaluation \(F\), with \(F=S\) or \(F=P\) in APP,
is \(r_3\circ F\). When reduction acts on the phase cursor, the same
balanced section produces the classes \(\{1,4,7\}\), \(\{2,5,8\}\) and
\(\{3,6,9\}\), with values \(+1,-1,0\). The TPK selector defined
next uses these classes to determine the operational mode.
Reading an evaluation and selecting a mode thus have the same
ternary reduction, but different domains and functions.

\Needspace{16\baselineskip}
\begin{proposition}[Arithmetic of balanced pairs]
\label{x_reciprocidad:trit-ampl-prop-aritmetica}
Let \(n=r+3c\) and \(m=s+3d\), with \(r,s\in\TRIT\). Define
\[
 \chi(r,s)=\frac{r+s-r_3(r+s)}3.
\]
The integer operations are transported to pairs through
\begin{align}
 (r,c)\boxplus(s,d)
 &=\bigl(r_3(r+s),c+d+\chi(r,s)\bigr),
 \label{x_reciprocidad:trit-ampl-eq-suma}\\
 (r,c)\boxtimes(s,d)
 &=\bigl(rs,rd+sc+3cd\bigr).
 \label{x_reciprocidad:trit-ampl-eq-producto}
\end{align}
The map \(\mathcal L\) preserves both operations.
\end{proposition}
\begin{proof}
The equality \(r+s=r_3(r+s)+3\chi(r,s)\) gives
\[
 n+m=r_3(r+s)+3\bigl(c+d+\chi(r,s)\bigr).
\]
On the other hand,
\[
 nm=rs+3(rd+sc+3cd).
\]
Since the product of two balanced representatives again belongs to
\(\{-1,0,+1\}\), it requires no additional residue reduction. Uniqueness
of the decomposition proves both formulas.
\end{proof}

For example, \(2=(-1)+3\) and \(4=1+3\). Their sum and product are reconstructed
as
\[
 (-1,1)\boxplus(1,1)=(0,2),\qquad
 (-1,1)\boxtimes(1,1)=(-1,3).
\]
The resulting evaluations are \(6\) and \(8\). Adding quotients suffices
in this additive example; the product involves the cross terms
\(rd\), \(sc\) and \(3cd\). The discrepancy between the two operations also lies
in the prolongation coordinate, not only in the residue.

The term \(\chi\) is a cocycle of the balanced section. If
\(r\oplus s=r_3(r+s)\), it satisfies
\[
 \chi(r,s)+\chi(r\oplus s,t)
 =\chi(s,t)+\chi(r,s\oplus t).
\]
Both sides represent the quotient of the same sum of three integers.
This identity guarantees that regrouping the operations does not change the
reconstructed evaluation. At the same time, projection onto the first
component suppresses that accounting: \(1+1=-1+3\), whereas the residual
reading preserves only \(-1\). The eliminated information is
determined, in this case, by \(\chi(1,1)=1\).

The relation with the modulus \(9\) preserves a second discrete scale.
The map
\[
 \beta:\TRIT\times\mathbb F_3\longrightarrow\mathbb Z/9\mathbb Z,
 \qquad (r,\bar c)\longmapsto r+3c\pmod9
\]
is well defined and bijective. Changing the representative of \(\bar c\)
adds a multiple of \(9\); conversely, equality of images first fixes
\(r\) by reduction modulo \(3\) and then \(\bar c\). The sum
transported by this bijection contains the term
\(\chi(r,s)\) of \eqref{x_reciprocidad:trit-ampl-eq-suma}, now reduced modulo \(3\).
Thus the two ternary levels are not composed through an independent
sum of coordinates: normalization of the first level modifies
the second.

In particular, the classes \(3\), \(6\) and \(0\) of
\(\mathbb Z/9\mathbb Z\) have first component \(0\), but second
component \(1\), \(2\) and \(0\), respectively. The distinction already observed
between direct reduction and extraction of the ideal coordinate appears
as a property of \(\beta\). Preserving that coordinate prevents the
zero residual sector from being artificially turned into a single position.
The modulus \(1000\), used later to organize decimal blocks,
serves another positional representation function; its quotient is preserved
through the corresponding division. A change of base preserves the integer
when it retains residue and quotient, whereas the coincidence of residues
in different bases does not by itself identify their transport states.

\Needspace{11\baselineskip}
\subsubsection{Reversal, residual class and nonvisible coordinates}
\label{x_reciprocidad:trit-ampl-inversion}

Additive inversion lifts unambiguously:
\[
 \mathcal B(-n)=(-r_3(n),-q_3(n)).
\]
In particular, \(3\), \(0\) and \(-3\) have coordinates
\((0,1)\), \((0,0)\) and \((0,-1)\). They share a zero signature, but represent
different evaluations. In a transition restricted to a quotient of
maximum amplitude \(1\), they are the three possibilities of the visible fiber over
zero; in a general integer lift, that fiber contains all pairs
\((0,c)\), with \(c\in\mathbb Z\).

The quotient coordinate recovers the evaluated integer, but does not by
itself recover the trajectory that produced it. Two APP paths may share
evaluation, residue and quotient, and differ in orientation, mode sequence or
winding. Thus the decomposition \(\mathcal B\) is integrated into the
transport state, instead of replacing it. The arithmetic memory of a
division and the memory of a trajectory are related coordinates, not
synonyms.

For a signature sequence of length \(n\), oriented reversal is
\[
 C_{\rm rev}(\tau_1,\ldots,\tau_n)
   =(-\tau_n,\ldots,-\tau_1).
\]
Applying it twice restores the sequence. The operation reverses the order of
the positions and the sign of each signature; it preserves its length and the
zero positions in reverse order. In the adopted parametrization,
it interchanges the return with memory, associated with \(+1\), and bidirectional
opening, associated with \(-1\), maintaining the threshold \(0\). The transport
of the remaining data of a trajectory is specified through the
corresponding involution of TPK.

This reversal of signatures must be distinguished from conjugation within a
quadratic algebra. The former may change the type \(\tau\); the latter,
which is constructed next, acts within the type already fixed. The
distinction prevents a sign reversal from being credited with an algebraic
equivalence between different regimes.

\subsubsection{Quadratic product, conjugation, and algebraic norm}
\label{x_reciprocidad:trit-ampl-norma}

An element of \(\mathbb A_\tau\) has a unique expression \(aI+bJ_\tau\).
The relation \(J_\tau^2=-\tau I\) determines its product:
\begin{equation}
 (aI+bJ_\tau)(cI+dJ_\tau)
   =(ac-\tau bd)I+(ad+bc)J_\tau.
 \label{x_reciprocidad:trit-ampl-producto-cuadratico}
\end{equation}
This is a real realization of the type already selected. The coordinates
\(a,b\) are not residues of APP, and the product
\eqref{x_reciprocidad:trit-ampl-producto-cuadratico} does not replace
\eqref{x_reciprocidad:trit-ampl-eq-producto}; it expresses another operation, in the declared
quadratic codomain.

\begin{definition}[Conjugation and algebraic norm]
\label{x_reciprocidad:trit-ampl-def-norma}
Conjugation of type \(\tau\) and its algebraic norm are
\[
 \overline{aI+bJ_\tau}=aI-bJ_\tau,\qquad
 N_\tau(aI+bJ_\tau)=a^2+\tau b^2.
\]
\end{definition}

\begin{proposition}[Multiplicativity and invertibility]
\label{x_reciprocidad:trit-ampl-prop-norma}
For \(z,w\in\mathbb A_\tau\),
\[
 z\bar z=N_\tau(z)I,\qquad
 N_\tau(zw)=N_\tau(z)N_\tau(w).
\]
Moreover, \(z\) is invertible if and only if \(N_\tau(z)\ne0\), and in that case
\[
 z^{-1}=\frac{\bar z}{N_\tau(z)}.
\]
\end{proposition}
\begin{proof}
The product of \(aI+bJ_\tau\) and \(aI-bJ_\tau\) is
\((a^2+\tau b^2)I\). Conjugation is multiplicative and the algebra is
commutative, so
\((zw)\overline{zw}=(z\bar z)(w\bar w)\). If the norm does not vanish, the
indicated formula provides the inverse. Conversely, if \(zw=I\),
multiplicativity implies \(N_\tau(z)N_\tau(w)=1\).
\end{proof}

The quadratic form \(N_\tau\) is positive definite only in the elliptic type.
For \(\tau=0\), it equals \(a^2\), and the element \(J_0\) is nonzero although its
square and norm are zero. For \(\tau=-1\), it equals \(a^2-b^2\); the
nonzero elements \(I+J_{-1}\) and \(I-J_{-1}\) have zero norm and their
product vanishes. The algebraic norm here is a multiplicative quadratic
form: it does not presuppose a positive vector space norm.

The three types can be faithfully realized through
\[
 aI+bJ_\tau\longmapsto
 \begin{pmatrix}a&-\tau b\\b&a\end{pmatrix}.
\]
The determinant of this matrix is \(N_\tau(aI+bJ_\tau)\). Thus the
complex structure, the dual number structure and the split real structure
are represented with the same matrix expression.
Each preserves its law and its zero divisors. The form of signature
\((1,1)\) in the last case has a Lorentzian interpretation in dimension
\(1+1\); identifying a physical metric additionally requires a specific space,
field and transport rule.

\subsubsection{Multiplicative conservation and composition of evolutions}
\label{x_reciprocidad:trit-ampl-evoluciones}

The exponential already obtained belongs to the set of elements of norm one:
\[
 N_\tau\bigl(\exp(tJ_\tau)\bigr)=1.
\]
This conservation is maintained when composing parameters. Indeed, all the
factors are functions of the same generator, and
\[
 \exp(tJ_\tau)\exp(sJ_\tau)=\exp((t+s)J_\tau).
\]
Comparing the two components yields the addition laws
\begin{align}
 C_\tau(t+s)&=C_\tau(t)C_\tau(s)-\tau S_\tau(t)S_\tau(s),\\
 S_\tau(t+s)&=S_\tau(t)C_\tau(s)+C_\tau(t)S_\tau(s).
\end{align}
Conjugation changes \(t\) to \(-t\) and coincides with inversion within
this subgroup. It does not change \(\tau\): it reverses an evolution of the chosen
regime.

In the elliptic type, the preserved form is \(a^2+b^2\), and the exponential
trajectory traverses the circle of norm one. In the hyperbolic type, the
eigencoordinates are \(u=a+b\), \(v=a-b\), with \(uv=N_{-1}(z)\).
The exponential gives
\[
 (u,v)=(e^t,e^{-t}),\qquad uv=1.
\]
The growth of one coordinate is accompanied by the contraction of the other.
Conservation expresses a relation between both coordinates; it does not state that
each one remains constant.

In the parabolic type,
\[
 (I+tJ_0)(I+sJ_0)=I+(t+s)J_0,\qquad
 (I+tJ_0)^{-1}=I-tJ_0.
\]
Evolution does not stop because \(J_0^2=0\). Nilpotency eliminates the
terms of order at least two, but preserves the linear term that
transports the parameter. The visible symbol \(0\), the
zero element of the algebra and the nonzero generator \(J_0\) are distinct. This distinction
is necessary to describe the threshold without identifying it with an absence of
state or dynamics.

\begin{proposition}[Composition of affine coordinates]
\label{x_reciprocidad:trit-ampl-prop-pendientes}
In the chart where the scalar component does not vanish, represent a
direction by \(I+uJ_\tau\). If \(1-\tau uv\ne0\), the direction of its
product with \(I+vJ_\tau\) has coordinate
\[
 u\star_\tau v=\frac{u+v}{1-\tau uv}.
\]
\end{proposition}
\begin{proof}
We have
\[
 (I+uJ_\tau)(I+vJ_\tau)
   =(1-\tau uv)I+(u+v)J_\tau.
\]
Dividing by the scalar component gives the stated expression.
\end{proof}

The affine chart preserves the ratio between components, not the scalar factor
divided out. When \(1-\tau uv=0\), the product must be examined before
changing charts: it may have a defined direction with zero scalar
component, or be the zero element in the presence of zero divisors. The composition
law still belongs to the full algebra; the singularity of
a coordinate does not automatically amount to a singularity of the object.
This formula will reappear in the elliptic regional reading. Its use requires
retaining the denominator and, where appropriate, the normalization factor.

\subsubsection{Transport of type and subsequent realizations}
\label{x_reciprocidad:trit-ampl-puentes}

The algebras \(\mathbb A_{+1}\), \(\mathbb A_0\) and
\(\mathbb A_{-1}\) are not isomorphic as unital real algebras.
The first is a field; the second contains a nonzero nilpotent; the
third has nontrivial idempotents and is isomorphic to
\(\mathbb R\times\mathbb R\). Thus the change of signature that organizes a
route cannot be interpreted as an undifferentiated isomorphism between the
three realizations. TPK preserves the regime index, the domains and the
effective operator acting at each transition.

Within the electronic chapter, the pair of projections of the APP sheets
produces a normalized commutator whose square is \(-I\). The same block
contains an involution of square \(I\) and two nilpotent directions. The
joint realization will be proved on that specific pair of projections;
the preceding classification explains why the three quadratic relations
can appear in the same matrix algebra without being identified with one another.

In the exceptional prolongation, the Paley matrix reduced modulo \(3\)
satisfies \(A_W^2=-I_6\). The relation endows the ternary space with a
structure of dimension \(3\) over \(\mathbb F_9\). Its link with a real
operator root is expressed through the integral presentation
\[
 \mathbb Z[u]/(u^2+1).
\]
The base changes to \(\mathbb F_3\) and to \(\mathbb R\) produce
representations of that same relation. Each preserves its characteristic,
its dimension and its incidence information. The development of the
respective operators establishes the link, not the equality of their
cardinalities.

The conservation of norm one therefore describes the quadratic
realizations of evolutions. The conservation of quotients, orientation and
trajectory belongs to the transported arithmetic state. Both properties
play a role in the construction, with different functions: one determines the
identities of the representations; the other allows the operations that originate them
to be composed, reconstructed and prolonged.

% END OWNER



% BEGIN OWNER /Users/ruben/Documents/New project/output/EXTREMA_MEDIA_RAZON_RECIPROCIDAD_ES_EN_20260918/ARTICULO/ES/source/figures/regimenes_trit.tex
\begin{figure}[htbp]
\centering
\begin{tikzpicture}[scale=.88,>=Stealth,font=\small]
\foreach \c/\titulo in {0/{Elliptic},4.6/{Parabolic},9.2/{Hyperbolic}}{
\begin{scope}[xshift=\c cm]
\draw[->,gray] (-1.65,0)--(1.7,0) node[right] {\(u\)};
\draw[->,gray] (0,-1.35)--(0,1.5) node[above] {\(v\)};
\node at (0,2.05) {\titulo};
\end{scope}}
\draw[thick] (0,0) circle[radius=1];
\draw[->,thick] (1,0) arc[start angle=0,end angle=60,radius=1];
\node at (0,-1.85) {\(u^2+v^2=1\)};
\node at (0,-2.35) {\(J^2=-I\)};
\draw[thick] (3.6,-1.2)--(3.6,1.2);
\draw[->,thick] (5.6,-1.2)--(5.6,1.2);
\node at (4.6,-1.85) {\(u^2=1\)};
\node at (4.6,-2.35) {\(J^2=0\)};
\draw[thick,domain=-.9:.9,samples=50] plot ({9.2+cosh(\x)},{sinh(\x)});
\draw[thick,domain=-.9:.9,samples=50] plot ({9.2-cosh(\x)},{sinh(\x)});
\node at (9.2,-1.85) {\(u^2-v^2=1\)};
\node at (9.2,-2.35) {\(J^2=I\)};
\end{tikzpicture}
\caption{Quadratic realizations of TRIT. For
\(J_\tau=\left(\begin{smallmatrix}0&-\tau\\1&0\end{smallmatrix}\right)\),
the exponential preserves \(u^2+\tau v^2\). Parabolic degeneration
preserves \(u\) and transports \(v\) linearly. The three diagrams represent
orbits of quadratic forms; their geometric type corresponds to the algebraic
relation of the generator.}
\label{x_reciprocidad:fig:trit-regimenes}
\end{figure}

% END OWNER


\newpage
\subsection{TPK: topological phase kernel}

The dynamics uses the phase selector
\[
 M_{\rm ph}(r)=
 \begin{cases}
 +1,&r\in\{1,4,7\},\\
 -1,&r\in\{2,5,8\},\\
 0,&r\in\{3,6,9\}.
 \end{cases}
\]
Its value decides which evaluation operates. Cardinal transport shifts the corresponding cursor; the register preserves the preceding value, the new quotient, the orientation and the transition. Therefore, selection and transport are distinct functions: a ternary sign is not a displacement of the torus.

We shall denote by \(\widehat X\) the space of arithmetic states with trajectory memory. An element contains the positions and directions of the two cursors, the active mode, the TRIT signature, the phase, the residues and quotients of both evaluations, the carries, the constructed prefix and the boundary data. When a refinement cylinder is introduced, it forms part of the state. The notation avoids replacing the object with a reduced version that preserves only the visible symbol.

\begin{definition}[TPK update]
Let \(\mathsf{Sel}_t\) be mode selection through \(M_{\rm ph}\), \(\mathsf{Tra}_t\) lifted cardinal transport and \(\mathsf{Mem}_t\) the update of the registers. The transition is
\begin{equation}
 \mathcal U_t=\mathsf{Mem}_t\circ\mathsf{Tra}_t\circ\mathsf{Sel}_t.
 \label{x_reciprocidad:eq:actualizacion}
\end{equation}
Each map preserves the fields that it does not modify and updates the others through Euclidean division and concatenation of the trajectory.
\end{definition}

The notation \eqref{x_reciprocidad:eq:actualizacion} expresses the order of composition; it does not replace the effective rules. The following section fully specifies its finite two-cursor realization. Afterward, prolongation will not be obtained by erasing its memory and repeating the first emission, but by transporting the prefixes and their compatible boundaries. In this way, APP provides the evaluations, TRIT distinguishes their regimes and TPK composes their evolution.


% BEGIN OWNER /Users/ruben/Documents/New project/output/EXTREMA_MEDIA_RAZON_RECIPROCIDAD_ES_EN_20260918/ARTICULO/ES/source/sections/tpk_desarrollo_integrado.tex
% Recuperación expositiva desde los propietarios del núcleo TPK.
% La procedencia y el control de conservación se documentan en technical/TPK_INTEGRACION.md.
% Este archivo amplía la sección TPK; no sustituye extension.tex ni generacion.tex.

\subsubsection{Dynamic state, observables and dependence of the operations}
\label{x_reciprocidad:tpk-int:estado}

The topological phase kernel organizes the evolution of the two APP sheets
with the orientation determined by TRIT. Its definition includes the state
on which it acts, the transition rules and the relations that make
the prolongations compatible. The equation \eqref{x_reciprocidad:eq:actualizacion} acquires
operational content by specifying which coordinates each
of its factors receives and modifies. Position belongs to the toroidal support; phase
belongs to the calendar; orientation determines transport; the
quotients record the arithmetic evaluation before its reduction.
These coordinates are related through the transition and preserve different
mathematical functions.

The observable projection of a state is
\begin{equation}
 q=(x_+,d_+,x_\times,d_\times,\theta)
 \in\mathcal Q_{\rm obs}
 :=(X_9\times\mathcal D)^2\times\Theta_{108},
 \qquad \mathcal D=\{N,E,S,O\},\quad
 \Theta_{108}=\ZZ/108\ZZ.
 \label{x_reciprocidad:tpk-int:qobs}
\end{equation}
The subscripts distinguish the additive and multiplicative cursors. In the
fibers of complete states over \(q\), the ordered trajectory, the integer evaluations
of both sheets, their residues and quotients, and the orientation
and boundary registers are preserved. In a precision prolongation, the
prefix and the compatible cylinder are added. Therefore, two states with the same
projection \(q\) may preserve different registers.

The dependence of the operations has three levels. At the local level,
the selector activates a sheet, transport modifies its position and the
update calculates the new arithmetic coordinates. At the trajectory
level, these transitions are composed and preserve the registers of
the preceding steps. At the refinement level, extension
relations determine which continuation preserves the prefix, the signature and the
boundary of the antecedent. The finite calendar describes the phase of these
operations; holonomy describes their composition on the history.

\subsubsection{Tritic selection and lifted cardinal transport}
\label{x_reciprocidad:tpk-int:seleccion}

Let \(\bar\theta\in\{0,\ldots,107\}\) be the calendar representative.
The operational phase and the dodecaphasic sector are
\begin{equation}
 r(\theta)=1+(\bar\theta\bmod9),\qquad
 \beta(\theta)=\left[\left\lfloor\frac{\bar\theta}{9}\right\rfloor\right]_{12}.
 \label{x_reciprocidad:tpk-int:fase}
\end{equation}
The selector already defined can be written \(\varepsilon_9=M_{\rm ph}\).
Its vector of values, in the ordered chart of \(D_9\), is
\[
 m_{\rm ph}=(1,-1,0,1,-1,0,1,-1,0).
\]
The function \(M_{\rm ph}:D_9\to\TRIT\) and this vector are two
presentations of the same selector, with their respective domains.

Define the activation indicators
\[
 a_+(\theta)=\mathbf1_{\{1,4,7\}}(r(\theta)),\qquad
 a_\times(\theta)=\mathbf1_{\{2,5,8\}}(r(\theta)).
\]
The transport of the cursors is then
\begin{equation}
 x_+'=x_++a_+(\theta)\nu_{d_+},\qquad
 x_\times'=x_\times+a_\times(\theta)\nu_{d_\times}
 \quad\text{in }X_9.
 \label{x_reciprocidad:tpk-int:cursores}
\end{equation}
In phases \(3,6,9\) both positions remain unchanged; the neutral event
is part of the chronology and its register. The cardinal directions
determine the vectors \(\nu_d\), while the phase signature determines
which of those vectors is applied.

The integer lift preserves the boundary crossing. If
\(s:X_9\to D_9^2\) chooses the positive representatives and
\(x'=x+a\nu_d\), with \(a\in\{0,1\}\), one defines
\begin{equation}
 b_d(x,a)=\frac{s(x)+a\nu_d-s(x')}{9}\in\ZZ^2.
 \label{x_reciprocidad:tpk-int:borde}
\end{equation}
Divisibility by \(9\) follows from equality in \(X_9\).
If \(z\in\ZZ^2\) records the preceding crossings, the update is
\(z'=z+b_d(x,a)\). In this way,
\[
 s(x')+9z'=s(x)+9z+a\nu_d.
\]
The integer transport datum remains reconstructible after reducing
the position on the torus.

\begin{proposition}[Conservation of the lifted displacement]
\label{x_reciprocidad:tpk-int:prop-desplazamiento}
For a cardinal trajectory of \(n\) steps, whose activation
indicators are \(a_t\), one has
\[
 s(x_n)+9z_n=s(x_0)+9z_0+
 \sum_{t=0}^{n-1}a_t\nu_{d_t}.
\]
The same equality is preserved when concatenating two compatible trajectories.
\end{proposition}
\begin{proof}
The one-step identity is \eqref{x_reciprocidad:tpk-int:borde}. Summing it over the
\(n\) steps cancels the intermediate coordinates. Dividing the summation
interval into two segments produces exactly the concatenation law. In a
closed trajectory, the difference \(z_n-z_0\) is the winding
number defined previously.
\end{proof}

The same principle governs APP evaluations. If a sheet takes the
integer value \(N_t\), its register simultaneously preserves
\(N_t=R_t+9q_t\). The difference between two steps satisfies
\[
 N_{t+1}-N_t=(R_{t+1}-R_t)+9(q_{t+1}-q_t).
\]
The spatial crossing memory \(z_t\) and the arithmetic quotient \(q_t\)
have different definitions; their joint conservation maintains the
geometric and arithmetic provenance of each transition.

\subsubsection{Two-cursor chronology and composite returns}
\label{x_reciprocidad:tpk-int:cronologia}

The cardinal involution \(\jmath\) interchanges \(N\) with \(S\) and
\(E\) with \(O\). After applying \eqref{x_reciprocidad:tpk-int:cursores}, both
directions are reversed when a block of \(27\) steps ends.
Writing
\[
 h(\theta)=\mathbf1_{\{0,27,54,81\}}
             ((\bar\theta+1)\bmod108),
\]
the complete rule on the observable projection is
\begin{equation}
 F_{\rm obs}(q)=
 \bigl(x_+',\jmath^{h(\theta)}d_+,
       x_\times',\jmath^{h(\theta)}d_\times,
       \theta+1\bigr).
 \label{x_reciprocidad:tpk-int:fobs}
\end{equation}
The reading that forms the emission uses the positions before
transport, according to the explicit recurrence in
section~\ref{x_reciprocidad:sec:extension}. The observable update and the emission
are therefore composed in a fixed order.

\begin{proposition}[Reversibility and period of the observable calendar]
\label{x_reciprocidad:tpk-int:periodo}
The map \(F_{\rm obs}\) is a permutation of
\(\mathcal Q_{\rm obs}\) and satisfies \(F_{\rm obs}^{108}=I\).
In a block of \(54\) steps, positions and directions are restored;
the coordinate of \(\Theta_{108}\) advances by \(54\).
\end{proposition}
\begin{proof}
Given the final state, the preceding phase is first recovered by subtracting one
in \(\Theta_{108}\). That phase determines whether the
directions were reversed and which cursor moved. Applying the corresponding involution
and subtracting the cardinal vector uniquely recovers the initial state.

Each interval of \(27\) steps contains nine activations of each
cursor. The next block preserves the activations and reverses the
directions, so the integer displacements cancel in
\(54\) steps. The rule has period \(54\) in its position
and orientation increments; the same cancellation holds with any origin
of the calendar. After \(108\) steps, \(\theta\) is also restored.
\end{proof}

We reserve
\[
 \mathcal K_{27}=F_{\rm obs}^{27}:
 \mathcal Q_{\rm obs}\longrightarrow\mathcal Q_{\rm obs}
\]
for the iterate of \(27\) transitions. Its fourth iterate is the
identity on the observable state defined in \eqref{x_reciprocidad:tpk-int:qobs}.
Projection of the history onto the calendar allows this
finite check to be performed. The trajectory register, its evaluations and
the new prefixes are preserved in the dynamic state on which
\(\mathcal U_t\) acts; observable equality does not reset them.

\subsubsection{Affine transport of memory and composition of states}
\label{x_reciprocidad:tpk-int:memoria}

The observable projection admits fibers of arithmetic and geometric data.
Over a configuration \(q\), we write a state in the form
\[
 x=(q,o,h,c,\sigma,C,b,\lambda).
\]
Here \(o\) preserves the orientation and the sheet; \(h\), the residues
and quotients; \(c\), the additive transport memory; \(\sigma\),
the boundary signature; \(C\), the precision cylinder; \(b\), its
boundary label; and \(\lambda\), the recorded trajectory. The
signature is determined by a map
\[
 \operatorname{Sig}_q:
 \mathsf H_q\times\mathsf C_q\times\mathsf{Hist}_q
 \longrightarrow\mathsf S_q,
 \qquad \sigma=\operatorname{Sig}_q(h,c,\lambda).
\]
The symbols \(\mathsf H_q,\mathsf C_q,\mathsf{Hist}_q\) and
\(\mathsf S_q\) denote, respectively, the sets of residue--quotient
data, the abelian group of transport memory, the
recorded trajectories and the boundary signatures.
The realizations of this signature are specified with their components in
the regional and terminal constructions. This dependence prevents treating
the signature as a free parameter once the data producing it are fixed.

The elementary residue and quotient chart is explicit. For
\(n\in\ZZ\), let
\[
 r=1+((n-1)\bmod9),\qquad u=\frac{n-r}{9}.
\]
Then \(n=r+9u\), with \(r\in D_9\) and \(u\in\ZZ\).
Uniqueness follows because two representatives of \(D_9\) congruent
modulo \(9\) are equal. This chart provides local integer
reconstruction. The precision towers additionally incorporate the truncation
maps and the cylinders developed in
\ref{x_reciprocidad:sec:extension} and \ref{x_reciprocidad:sec:generacion}.

Let \(\gamma:q\to q'\) be an observable trajectory. The transport
of the fibers is expressed through maps
\[
 \mathsf H(\gamma):\mathsf H_q\to\mathsf H_{q'},\qquad
 \mathsf C(\gamma):\mathsf C_q\to\mathsf C_{q'},\qquad
 \mathsf{Hist}(\gamma):\mathsf{Hist}_q\to\mathsf{Hist}_{q'}.
\]
The maps respect identities and composition, and
\(\mathsf C(\gamma)\) is a homomorphism of abelian groups.
The increment \(k_{\rm tr}(\gamma)\in\mathsf C_{q'}\) satisfies
\begin{equation}
 \begin{split}
 k_{\rm tr}(\operatorname{id}_q)&=0,\\
 k_{\rm tr}(\gamma_2\circ\gamma_1)
 &=k_{\rm tr}(\gamma_2)+
   \mathsf C(\gamma_2)k_{\rm tr}(\gamma_1).
 \end{split}
 \label{x_reciprocidad:tpk-int:cociclo}
\end{equation}
The symbol \(k_{\rm tr}\) here designates the transport cocycle;
the rational coordinate \(\kappa\) of the register \(K\) retains
its notation in the chapter on dodecaphasic closure.

The transportable coordinates are updated by
\begin{equation}
 \begin{aligned}
 h'&=\mathsf H(\gamma)h,&
 c'&=\mathsf C(\gamma)c+k_{\rm tr}(\gamma),\\
 \lambda'&=\gamma\circ\lambda,&
 \sigma'&=\operatorname{Sig}_{q'}(h',c',\lambda').
 \end{aligned}
 \label{x_reciprocidad:tpk-int:accion-memoria}
\end{equation}
The second equation is an affine action: the linear term transports
the preceding memory and the cocycle adds the increment produced by the
trajectory. In the realization of spatial crossing for each cursor separately,
\(\mathsf C(\gamma)=I\) and \(k_{\rm tr}(\gamma)\) is the sum
of the vectors of \eqref{x_reciprocidad:tpk-int:borde}. This concrete instance
links the abstract law to the cardinal operations already defined.

\begin{proposition}[Compatibility of the update with concatenation]
\label{x_reciprocidad:tpk-int:composicion-memoria}
The update \eqref{x_reciprocidad:tpk-int:accion-memoria} by
\(\gamma_1\), followed by the update by \(\gamma_2\),
coincides with the update by \(\gamma_2\circ\gamma_1\).
\end{proposition}
\begin{proof}
For the affine coordinate, two updates give
\[
 c''=\mathsf C(\gamma_2)\mathsf C(\gamma_1)c
       +\mathsf C(\gamma_2)k_{\rm tr}(\gamma_1)
       +k_{\rm tr}(\gamma_2).
\]
The functoriality of \(\mathsf C\) and
\eqref{x_reciprocidad:tpk-int:cociclo} turns this expression into
\(\mathsf C(\gamma_2\circ\gamma_1)c+
k_{\rm tr}(\gamma_2\circ\gamma_1)\).
The assertion for \(h\) follows from the composition of
\(\mathsf H\), and for \(\lambda\), from the associativity of the
trajectories. The final signature coincides because it is evaluated on the same
three reconstructed coordinates.
\end{proof}

The cylinder and its boundary complete this law through the relations
\(\operatorname{Ext}_\gamma\), defined in
\ref{x_reciprocidad:tpk-int:bidireccionalidad}. Their elements determine which of
the preceding updates preserve an admissible prolongation.
The states, together with the triples \((\gamma,x,x')\) admitted by
these relations, form a category over the observable trajectories:
\[
 (\gamma_2,x',x'')\circ(\gamma_1,x,x')
     =(\gamma_2\circ\gamma_1,x,x'').
\]
The identity of \(x\) is \((\operatorname{id}_q,x,x)\).
Composition preserves both arithmetic transport and boundary
prolongability. This is the structure in which returns
with memory and successive refinements are carried out.

\subsubsection{Dodecaphasic register, phase and accumulated memory}
\label{x_reciprocidad:tpk-int:dodecafase}

The \(108\) transitions are distributed in the ordered windows
\[
 E_m=\{9m-8,\ldots,9m\},\qquad 1\leq m\leq12.
\]
The register preserves the phase origin, the orientation and the data
produced during each window. Its map is
\[
 \mathcal R_{12}(x)=
 \bigl((E_m,\tau_m,\sigma_m,\kappa_m,\Lambda_m)\bigr)_{m=1}^{12},
\]
where \(\tau_m\) is the subroute of the window, \(\sigma_m\) its
orientation, \(\kappa_m\) the integer boundary increment and
\(\Lambda_m\) the local incidence register. The chart
that will be developed in \ref{x_reciprocidad:subsec:k-registro-integral} is adopted; its window
indices correspond to \(m=\beta+1\).
The domain comprises the compatible histories
of TPK and the codomain, their oriented temporal registers. The
composition of segments and their increments is governed by
\eqref{x_reciprocidad:tpk-int:accion-memoria}. The group \(C_{12}\) instead describes
the translation of the window index; \(\Theta_{108}\) preserves
the step calendar. This distinction makes it possible to specify what is transformed
and what is observed upon completing a cycle.

In the coordinates \(\theta=9b+r\), with \(0\leq r<9\), advancing
by \(9\) steps preserves \(r\) and increases \(b\) by one.
The addition of phases introduces the quotient
\[
 c_0(r,r')=\left\lfloor\frac{r+r'}9\right\rfloor,
\]
whose cocycle identity ensures associativity of composition.
The finite calendar preserves \([b]_{12}\); a turn register
preserves \(b\in\ZZ\), or \(b\in\mathbb N_0\) for
a history with fixed origin. Section~\ref{x_reciprocidad:sec:extension} proves
the nonsplit exact sequence \eqref{x_reciprocidad:eq:extension108}. Its function is
to describe the relation between these coordinates, before using them in the
construction of the signed state.

The signed register receives the phase and memory data through the
composition developed in section~\ref{x_reciprocidad:sec:k-registro}:
\begin{equation}
 R_{\rm sgn}=\Pi_H\circ\mathcal E_{108}^{90,120}
                         \circ\mathcal R_{12}.
 \label{x_reciprocidad:tpk-int:registro-causal}
\end{equation}
Its domain is the terminal state with memory; its output belongs to
\(\ZZ^{12}\). The reversible transformation between that signed state
and \(K\), as well as the rational reading of \(K\), are proved in
the corresponding chapter. The register operations precede
these transformations. Reconstructing a register from its
cyclic differences verifies its recoverability; dynamic provenance
must be preserved through the events that produced it.

For the encoded trace \((d_t)\) preserved by the register, the
event reader counts the codes \(\{2,4,6,8\}\) of each
window with the orientation
\((+,+,+,-,-,-,+,+,+,-,-,-)\). These codes and the cardinal
symbols of \(\mathcal D\) are distinct coordinates of the
register. Chapter~\ref{x_reciprocidad:sec:k-registro} develops this reader,
the integral decomposition \(1+5+6\), the congruences that
characterize its image and its exact inverse. The inversion proof
establishes which data that chart preserves; the definition of
\(\mathcal R_{12}\) specifies the oriented history on which
it acts.

\subsubsection{Phase incidence and internal construction of the channels}
\label{x_reciprocidad:tpk-int:incidencia}

The selector determines three subsets of \(D_9\):
\[
 H_+=\{1,4,7\},\qquad H_-=\{2,5,8\},\qquad H_0=\{3,6,9\}.
\]
Let \(H_{\rm act}=H_+\sqcup H_-\) be the set of active phases and
let \(O_{\rm ph}=D_9\setminus\{9\}\) be the chart before the marked
return. Denote by \(P_2(H_{\rm act})\) the set of
unordered pairs of active phases. The internal incidence is given by
\begin{equation}
 \mathcal F^{\rm ph}_6=H_{\rm act}\times P_2(H_{\rm act}),\qquad
 \mathcal F^{\rm ph}_8=O_{\rm ph}\times P_2(H_{\rm act}),\qquad
 \Theta_{\rm ph}=D_9\times P_2(H_{\rm act}).
 \label{x_reciprocidad:tpk-int:fases90120}
\end{equation}
The indices describe the number of positions of the first component.
The construction uses the phases and the calendar origin already fixed;
subsequent exceptional realizations will transport these
incidences to their own domains.

\begin{proposition}[Cardinalities and oriented phase incidence]
\label{x_reciprocidad:tpk-int:cardinales}
The sets in \eqref{x_reciprocidad:tpk-int:fases90120} have cardinalities
\(90,120,135\), respectively. If
\[
 \partial_{\rm ph}=\Theta_{\rm ph}\times\{-1,+1\},\qquad
 E_{\rm ph}=\partial_{\rm ph}\sqcup\{\infty\},\qquad
 L_{\rm ph}=H_0\times P_2(H_{\rm act}),
\]
then
\[
 |\partial_{\rm ph}|=270,\quad |E_{\rm ph}|=271,\quad
 |L_{\rm ph}|=45,\quad
 |\partial_{\rm ph}\sqcup L_{\rm ph}|=315.
\]
Simultaneously translating the phase origin and the marked return preserves
all these cardinalities and their inclusions.
\end{proposition}
\begin{proof}
There are six active phases, so
\(|P_2(H_{\rm act})|=\binom62=15\). The products by
\(6,8,9\) give \(90,120,135\). Orientation doubles \(135\);
the return mark adds one element; the three neutral phases give
\(3\cdot15=45\). A calendar translation transports each
factor through a bijection and takes the marked point to its image. The
products and disjoint unions thus preserve incidence.
\end{proof}

The cardinality \(271\) here has a realization through phase
incidence. The cubic completion of section~\ref{x_reciprocidad:sec:extension} has
its own construction of the crown \(1000-729\). The correspondence
between both realizations must preserve the maps and marks that
identify them, in addition to the cardinality. Likewise, the selector
\(\varepsilon_9\), the incidence \((90,120)\) and the extractor
\(\mathcal E_{108}^{90,120}\) preserve the domains and operations
that have just been distinguished.

\subsubsection{Reversal of trajectories and conjugate prolongations}
\label{x_reciprocidad:tpk-int:bidireccionalidad}

Bidirectionality begins with an operation on trajectories, before
assigning scalar values to them. For
\(\gamma=(x_0;d_1,\ldots,d_n)\), with endpoint \(x_n\), we define
\[
 \operatorname{rev}\gamma
   =(x_n;\jmath d_n,\ldots,\jmath d_1).
\]
The origin of the reversed trajectory is the endpoint of the original.
Composition satisfies
\[
 \operatorname{rev}^2\gamma=\gamma,\qquad
 \operatorname{rev}(\gamma_2\circ\gamma_1)
   =\operatorname{rev}\gamma_1\circ\operatorname{rev}\gamma_2,
\]
and the integer displacement changes sign. These identities are obtained
by reversing the order of the steps and replacing each cardinal vector with
its opposite. For a loop it follows that
\(\operatorname{wind}(\operatorname{rev}\gamma)
=-\operatorname{wind}(\gamma)\).

On the registers, reversal also requires transporting the order
of their components and the orientation of their events. Route reversal
and a signature conjugation are distinguishable operations: their
composition is specified when applied to a concrete register. In
particular, the regional involution that interchanges the closure
and propagation components preserves the autoscale axis; its formula is
presented in section~\ref{x_reciprocidad:mono-ampl:regiones}. This regional operation
acts on a different domain from the reversal of a cardinal trajectory.

Extensions over a route \(r:q\to q'\) are expressed through
relations between fibers of compatible states. If \(\mathcal F_q\)
denotes the fiber of complete states over \(q\), one requires
\[
 \operatorname{Ext}_r\subseteq\mathcal F_q\times\mathcal F_{q'},
 \qquad \Delta_{\mathcal F_q}\subseteq
        \operatorname{Ext}_{\operatorname{id}_q},
\]
\[
 \operatorname{Ext}_{r_2}\circ\operatorname{Ext}_{r_1}
       \subseteq\operatorname{Ext}_{r_2\circ r_1}.
\]
Their composition preserves the evaluations of both sheets, the orientation,
the quotients, the prefix and the boundary. Reversal of the route does not
by itself imply that every extension relation is a bijection.
The conjugate prolongation is defined on the domain where the transport of
these data satisfies the corresponding compatibilities.

\subsubsection{Nonadic holonomy and continuity of refinement}
\label{x_reciprocidad:tpk-int:holonomia}

One phase turn composes the transitions on the state with memory:
\begin{equation}
 \Gamma_{9,k}=\mathcal U_{k+8}\circ\cdots\circ\mathcal U_k.
 \label{x_reciprocidad:tpk-int:gamma}
\end{equation}
On an individualized history, this functional composition realizes
the holonomy relation; on the complete domain one preserves
\(\Gamma_{9,k}\subseteq\widehat X_k\times\widehat X_{k+9}\),
with the admissible prolongations developed later.
The cyclic base has nine phases, while
the fiber preserves the trajectory, the quotients and the boundary. The
monodromy is the action of one turn of that base on the fibers; the
holonomy \(\Gamma_{9,k}\) expresses its realization through the effective
transitions of TPK.

If \(g\) observes the phase and \(M\) the number of preserved turns,
\[
 g(\Gamma_{9,k}x)=g(x),\qquad M(\Gamma_{9,k}x)=M(x)+1.
\]
The first identity expresses phase closure and the second identifies
the new temporal depth. Coinductive prolongation also preserves
its cylinders and prefixes. The restrictions between depths
compose an inverse system; a compatible section gathers its
antecedents without replacing them with the latest observation. The development
of this arithmetic of histories and its unity occupies
\ref{x_reciprocidad:hist:seccion}--\ref{x_reciprocidad:hist:doble-varianza}; the construction of the
regions and their readers is carried out in section~\ref{x_reciprocidad:sec:generacion}.

Phase transfer on emissions already constructed has another
domain. Let \(\mathcal A_+\) and \(\mathcal A_\times\) be the
alphabets of additive and multiplicative signatures, of cardinalities \(18\)
and \(26\). Their product identifies the catalog \(\mathcal U_6\).
For \(f:\mathcal A_+\times\mathcal A_\times\to\RR\),
the sheet averages are
\[
 (E_+f)(s,e)=\frac1{18}\sum_{s'\in\mathcal A_+}f(s',e),\qquad
 (E_\times f)(s,e)=\frac1{26}\sum_{e'\in\mathcal A_\times}f(s,e').
\]
These emissions and their cardinalities are constructed through
the recurrence in section~\ref{x_reciprocidad:sec:extension}. Once they have been obtained,
the transfer operator is defined on
\(\mathcal U_6\times C_9\) by
\[
 P_+=E_+,\quad P_-=E_\times,\quad P_0=I,\qquad
 \mathcal K_{\rm ph}((u,r),(v,r+1))=P_{M_{\rm ph}(r)}(u,v),
\]
with zero weight for the remaining phase transitions. The index
\(r\) uses the positive representative of the cyclic class.
The phase order \((+1,-1,0)^3\) then produces the renewal
\[
 \mathcal K_{\rm ph}^{9}\big|_r=E_+E_\times.
\]
Indeed, both averages are idempotent, commute and average different
factors; their composition assigns weight \(1/(18\cdot26)=1/468\)
to each emission. This equality characterizes a subsequent transfer
observable. The transport of histories remains given by
\eqref{x_reciprocidad:tpk-int:gamma}, with its memory and its boundary.

The resulting hierarchy establishes the function of the following chapters.
The finite emission constructs the domain and the regions; the extension
relations preserve their prefixes; the nonadic connection prolongs them; the
dodecaphasic register determines the data involved in \(K\) and
in the closure of \(\alpha\). The numerical realization and the incidence
realization receive that preceding structure. This order allows
each consequence to be presented in its chapter while preserving, from the formal kernel,
the objects that make it possible.

% END OWNER



% BEGIN OWNER /Users/ruben/Documents/New project/output/EXTREMA_MEDIA_RAZON_RECIPROCIDAD_ES_EN_20260918/ARTICULO/ES/source/sections/memoria_resolvente.tex
\subsection{Accumulated memory and transport resolvents}
\label{x_reciprocidad:subsec:memoria-resolvente}

The return of the calendar does not eliminate the increments produced during
the path. To express this difference in operator terms, one first constructs
an accumulative register of trajectory events and then its
generating series. The reduction of memory variables will be carried out on
that transport, also preserving the source and the output reader.

\subsubsection{Oriented events and register update}

Let \((d_t)\) be the sequence of integer direction codes preserved
by a TPK trajectory. The code \(d_t\) and the cardinal direction
of the cursor are distinct coordinates; no identification
between their alphabets is assumed. We number the windows from zero
\[
 I_m=\{9m+1,\ldots,9m+9\},\qquad m\geq0.
\]
During the first cycle, \(I_m\) is window \(E_{m+1}\) of the
dodecaphase record. Its orientation is
\(\sigma=(1,1,1,-1,-1,-1,1,1,1,-1,-1,-1)\).
For a prolonged history, \(\sigma_m\in\{-1,1\}\) is the sign
retained in the corresponding window. The signed event is
\begin{equation}
 a_m=\sigma_m\sum_{t\in I_m}
       \mathbf1_{\{2,4,6,8\}}(d_t),\qquad |a_m|\leq9.
 \label{x_reciprocidad:mem:evento}
\end{equation}
Each summand is incorporated when its transition occurs. This definition
counts the events of the trace; it does not choose directions through a value
that is to be obtained later.

Let \(\mathbf e_0,\ldots,\mathbf e_{11}\) be the basis of \(\ZZ^{12}\)
and \(S\mathbf e_j=\mathbf e_{j+1\bmod12}\). We write
\(R=S^{-1}\). The accumulated register in fixed coordinates satisfies
\[
z_{m+1}=z_m+a_m\mathbf e_{m\bmod12}.
\]
A prefix without accumulated history starts at \(z_0=0\); at another
origin, \(z_0\) preserves the previous register.
The frame accompanying the calendar is \(y_m=S^{-m}z_m\).

\begin{proposition}[Update and return with memory]
\label{x_reciprocidad:mem:prop-actualizacion}
The register in the moving frame satisfies
\begin{equation}
 y_{m+1}=Ry_m+\eta_m,\qquad
 \eta_m=a_mR\mathbf e_0.
 \label{x_reciprocidad:mem:actualizacion}
\end{equation}
For all \(n\geq1\),
\begin{equation}
 y_n=R^ny_0+\sum_{j=0}^{n-1}R^{n-1-j}\eta_j.
 \label{x_reciprocidad:mem:iteracion}
\end{equation}
In particular,
\begin{equation}
 y_{12}=y_0+\sum_{j=0}^{11}a_j\mathbf e_j.
 \label{x_reciprocidad:mem:retorno12}
\end{equation}
\end{proposition}
\begin{proof}
Multiplying the update of \(z_m\) by \(S^{-(m+1)}\) gives
\(y_{m+1}=S^{-1}y_m+a_mS^{-1}\mathbf e_0\), because
\(S^{-m}\mathbf e_{m\bmod12}=\mathbf e_0\).
Formula \eqref{x_reciprocidad:mem:iteracion} is obtained by induction.
For \(n=12\), \(R^{12}=I\) and
\(R^{11-j}\eta_j=a_jR^{12-j}\mathbf e_0=a_j\mathbf e_j\),
which proves \eqref{x_reciprocidad:mem:retorno12}.
\end{proof}

Thus twelve windows restore the frame and preserve the event vector.
The count alone does not recover the order of all steps within each
window; that order remains in the recorded trajectory.

To specify the balance, define
\[
 D_3=I-S^3,\quad D_4=I-S^4,\quad
 B=D_3^{\mathsf T}D_3+D_4^{\mathsf T}D_4,
\]
\[
 \mathcal M(y)=\tfrac12y^{\mathsf T}By,
 \qquad q(y)=\sum_{j=0}^{11}y_j.
\]
The first is a quadratic form of the register; the second records
the sum of its coordinates. They are not identified here with physical energy or charge.

\begin{proposition}[Balance produced by the event]
\label{x_reciprocidad:mem:prop-balance}
The update \eqref{x_reciprocidad:mem:actualizacion} satisfies
\begin{equation}
 \begin{aligned}
 \mathcal M(y_{m+1})-\mathcal M(y_m)
   &=a_m(By_m)_0+2a_m^2,\\
 q(y_{m+1})-q(y_m)&=a_m.
 \end{aligned}
 \label{x_reciprocidad:mem:balance}
\end{equation}
\end{proposition}
\begin{proof}
Expanding the products,
\(B=4I-S^3-S^{-3}-S^4-S^{-4}\); therefore
\(R^{\mathsf T}BR=B\) and \(B_{00}=4\).
Since \(y_{m+1}=R(y_m+a_m\mathbf e_0)\), the quadratic difference
is \(a_m\mathbf e_0^{\mathsf T}By_m+\tfrac12a_m^2B_{00}\).
The coordinate sum is invariant under the permutation \(R\),
which gives the second equality.
\end{proof}

\Needspace{20\baselineskip}
\subsubsection{Generating series and uniform control of memory}

The series preserves all increments, not only their sum at the end of
a cycle. With a formal variable \(u\), let
\[
 Y(u)=\sum_{m\geq0}y_mu^m,\qquad
 H_\eta(u)=\sum_{m\geq0}\eta_mu^m.
\]

\begin{proposition}[Resolvent and tail bound]
\label{x_reciprocidad:mem:prop-serie}
The formal identity
\begin{equation}
 Y(u)=(I-uR)^{-1}\bigl(y_0+uH_\eta(u)\bigr).
 \label{x_reciprocidad:mem:serie}
\end{equation}
holds. Both series converge for \(|u|<1\). If \(\ell:\RR^{12}\to\RR\)
is linear, \(L=\|\ell\|_{1\to\RR}\), \(M_0=\|y_0\|_1\), and
\(0<r<1\), then, for every integer \(N\geq0\) and \(|u|\leq r\),
\begin{equation}
 \left|\sum_{m>N}\ell(y_m)u^m\right|
 \leq Lr^{N+1}\left(
 \frac{M_0+9(N+1)}{1-r}+\frac{9r}{(1-r)^2}\right).
 \label{x_reciprocidad:mem:cota}
\end{equation}
The series of derivatives converges uniformly on each disk
\(|u|\leq r<1\).
\end{proposition}
\begin{proof}
Multiplying \eqref{x_reciprocidad:mem:actualizacion} by \(u^{m+1}\) and summing gives
\(Y-y_0=uRY+uH_\eta\). The constant term of \(I-uR\) is
invertible, so its formal inverse is \(\sum_{k\geq0}u^kR^k\).
Moreover, \(R\) preserves the \(\ell^1\) norm and
\(\|\eta_m\|_1\leq9\), whence
\(\|y_m\|_1\leq M_0+9m\).
For the tail, sum the majorants
\(L(M_0+9m)r^m\), using
\[
 \sum_{m>N}r^m=\frac{r^{N+1}}{1-r},\qquad
 \sum_{m>N}mr^m=r^{N+1}
 \left(\frac{N+1}{1-r}+\frac r{(1-r)^2}\right).
\]
Finally, \(\sum_{m\geq1}Lm(M_0+9m)r^{m-1}<\infty\)
majorizes the derivative series and proves its uniform convergence.
\end{proof}

The variable \(u\) counts windows of nine transitions. Its identification
with the variable of another reader requires preserving that clock and the map
relating the two readings.

\subsubsection{Memory elimination: operator, source and reader}

On a domain where the update is fixed as \(U:X\to X\),
the free vector space \(V=\QQ^{(X)}\) on the states admits the operator
\(T\delta_x=\delta_{U(x)}\). The clock is included in \(x\) if
the transition depends on it. This linear extension preserves distinct
states even when they share an observable phase.

Consider a decomposition \(V=V_0\oplus V_1\), where \(V_1\)
is the auxiliary sector to be eliminated, and write
\[
 T=\begin{pmatrix}A&B_1\\C&D\end{pmatrix},\qquad
 v=\binom{v_0}{v_1},\qquad \ell=(\ell_0,\ell_1).
\]
Here \(A:V_0\to V_0\), \(B_1:V_1\to V_0\),
\(C:V_0\to V_1\), and \(D:V_1\to V_1\).
The following identities are formal; they also hold wherever the
resolvents converge.

\begin{proposition}[Schur reduction with source and reader]
\label{x_reciprocidad:mem:prop-schur}
Let
\[
 D_u=(I-uD)^{-1},\qquad
 \mathscr F(u)=I-uA-u^2B_1D_uC.
\]
The retained component of \((I-uT)^{-1}v\) is
\begin{equation}
 w_0=\mathscr F(u)^{-1}\bigl(v_0+uB_1D_uv_1\bigr),
 \qquad w_1=D_u(v_1+uCw_0).
 \label{x_reciprocidad:mem:schur-fuente}
\end{equation}
Its complete reading satisfies
\begin{equation}
 \begin{split}
 \ell(I-uT)^{-1}v
 ={}&(\ell_0+u\ell_1D_uC)
       \mathscr F(u)^{-1}(v_0+uB_1D_uv_1)\\
    &+\ell_1D_uv_1.
 \end{split}
 \label{x_reciprocidad:mem:schur-lector}
\end{equation}
\end{proposition}
\begin{proof}
The two equations of \((I-uT)w=v\) are
\(w_0=v_0+uAw_0+uB_1w_1\) and
\(w_1=v_1+uCw_0+uDw_1\).
Solving the second and substituting it into the first gives
\eqref{x_reciprocidad:mem:schur-fuente}; applying \(\ell_0\) and \(\ell_1\)
to both components gives \eqref{x_reciprocidad:mem:schur-lector}.
All the formal inverses exist because their constant terms
are identities.
\end{proof}

The term \(u^2B_1D_uC=\sum_{k\geq0}u^{k+2}B_1D^kC\)
describes an exit to memory, \(k\) steps within it and
a return. By contrast, \(u\ell_1D_uC\) reads the memory before that
return. Thus reducing only the operator does not necessarily preserve
the observation: source and reader must also be transformed.

The difference is visible in the cycle already constructed. If \(P\) retains
\(\mathbf e_0\), then \(PRP=0\), but
\begin{equation}
 \left\langle\mathbf e_0,(I-uR)^{-1}\mathbf e_0\right\rangle
 =\frac1{1-u^{12}},\qquad
 \left\langle\mathbf e_{11},(I-uR)^{-1}\mathbf e_0\right\rangle
 =\frac u{1-u^{12}}.
 \label{x_reciprocidad:mem:ejemplo-ciclo}
\end{equation}
Indeed, \(R^n\mathbf e_0=\mathbf e_{-n\bmod12}\): the
two series collect, respectively, \(n\equiv0\) and
\(n\equiv1\pmod{12}\). The resolvent of the compression is \(I\),
whereas the compression of the resolvent preserves the omitted returns.
This example corresponds to the twelve-window register, not to the entirety
of the TPK state.

\subsubsection{Composition of segments and coherence of reduction}

Let three composable affine transports be
\(F_i(y)=R_i y+b_i\), with compatible domains and codomains.
The memory produced by the complete path is
\begin{equation}
 F_3F_2F_1(y)
 =R_3R_2R_1y+b_3+R_3b_2+R_3R_2b_1.
 \label{x_reciprocidad:mem:cociclo-tres}
\end{equation}
The proof is the successive substitution of the three affine formulas.
Grouping the first two segments or the last two first produces
the same expression; the factors transport each increment from
the point where it originated. In \eqref{x_reciprocidad:mem:actualizacion},
\(R_i=R\) and \(b_i\) is the event of its window.

Eliminating several layers preserves a similar coherence.
To make it explicit, now decompose
\(V=V_0\oplus V_1\oplus V_2\) and write
\[
 T=\begin{pmatrix}
 A&B_1&B_2\\C_1&D_{11}&D_{12}\\C_2&D_{21}&D_{22}
 \end{pmatrix},\qquad Q_2=(I-uD_{22})^{-1}.
\]
Eliminating the third component leaves the four blocks
\begin{equation}
 \begin{aligned}
 A'&=A+uB_2Q_2C_2,&
 B_1'&=B_1+uB_2Q_2D_{21},\\
 C_1'&=C_1+uD_{12}Q_2C_2,&
 D_{11}'&=D_{11}+uD_{12}Q_2D_{21}.
 \end{aligned}
 \label{x_reciprocidad:mem:schur-intermedio}
\end{equation}

\begin{proposition}[Transitivity of elimination]
If \(D_*=(D_{ij})_{i,j=1}^2\), the complement obtained in two stages
coincides with that obtained by eliminating both memories together:
\begin{equation}
 \begin{split}
 I-uA'-u^2B_1'(I-uD_{11}')^{-1}C_1'
 ={}&I-uA\\
 &-u^2(B_1\ B_2)(I-uD_*)^{-1}\binom{C_1}{C_2}.
 \end{split}
 \label{x_reciprocidad:mem:schur-transitivo}
\end{equation}
\end{proposition}
\begin{proof}
Solving the third equation of \((I-uT)w=v\) and substituting it into the
first two produces \eqref{x_reciprocidad:mem:schur-intermedio}. Then eliminating the
second component gives the left-hand side of
\eqref{x_reciprocidad:mem:schur-transitivo}. Solving the two memory equations simultaneously
gives the right-hand side, by Proposition~\ref{x_reciprocidad:mem:prop-schur}.
The formal solution is unique because \(I-uT\) has constant term
\(I\). Both procedures also transform the source and
the reader according to \eqref{x_reciprocidad:mem:schur-fuente}--\eqref{x_reciprocidad:mem:schur-lector}.
\end{proof}

The balance \eqref{x_reciprocidad:mem:balance} describes the increments of the register.
The preceding identities allow them to be transported when reusing memory
in the dodecaphasic closure; they do not by themselves determine additional
analytic coefficients.

% END OWNER

